\section{The complete odd-minus estimate from the character-length bound}
\label{sec:odd-minus-bounds}

\status{OB1--OB8 reconstruct the full quantitative supplier
\texttt{ThmBcFree} from OS1 and \texttt{PsiBound}, with the explicit
constant $240$.  The proof includes the exceptional primes, the
original-field and local-data comparisons, the negative cyclotomic
twist, and all finite prime and twist sums.  The bespoke dependencies
are the literal \texttt{PsiBound} proved in OS9 and the finite Jacobi
annihilation proved in OS1.  The classical inputs listed below are
standard theorems, not new project-specific hypotheses.  In particular,
\texttt{killPow}, \texttt{gen}, \texttt{prodVal}, \texttt{atPrime} and
\texttt{deg} are proved here in the scope needed by the endpoint, rather
than retained as names for unresolved arguments.}

Fix a CM number field $E$, its given algebraic closure $\Omega$, a
prime $\ell$, and the source's local-data record $D$.  Thus $D$ lists
exactly the primes of $E$ above $\ell$, chooses a place of $\Omega$
above each, and realizes the cyclotomic character on the quotient
$G_{E,S}$ unramified outside those primes.  Put
\[
 A_\ell(e)=(\mathbb Q_\ell/\mathbb Z_\ell)(-e),\qquad
 \mathcal S_\ell(E,e;D)
 =\ker\!\left(H^1(G_{E,S},A_\ell(e))
       \longrightarrow\prod_{v\mid\ell}H^1(E_v,A_\ell(e))\right).
 \tag{OB.1}
\]
All cochains are continuous and the coefficients are discrete.  This is
the literal \texttt{Sha D e}, since \texttt{qzTw D e} has action
$\chi_{\rm cyc}^{-e}$.  Complex conjugation acts on this cohomology by
conjugation on the group and the same negative twist on coefficients.
For $e\ge1$ write
\[
 \mathcal S_\ell^-(E,e;D)
 =\{x\in\mathcal S_\ell(E,e;D):cx=(-1)^{e+1}x\},\qquad
 m_\ell(E,e;D)=\operatorname{length}_{\mathbb Z_\ell}
                       \mathcal S_\ell^-(E,e;D).
 \tag{OB.2}
\]
Initially lengths take values in $\mathbb N\cup\{\infty\}$; we prove
finiteness before treating them as real numbers.  The notation agrees
with \texttt{ShaMinus} and \texttt{minusLen}.  Write $n=[E:\mathbb Q]$,
and let $N\ge1$ be any modulus equipped with a rational embedding
$E\hookrightarrow\mathbb Q(\zeta_N)$.

\subsection{Classical inputs and the exact remaining interface}
\label{sec:ob-classical}

The following classical theorems are used in the precise forms stated.
They are recorded separately so that their use cannot be mistaken for
a reconstruction of a deeper project-specific supplier.
\begin{description}[style=nextline,leftmargin=0pt,labelwidth=0pt,labelsep=0pt]
\item[OB-CF: unramified class-field reciprocity.]
\label{ob-ref-classfield}
For a totally complex number field $K$, its maximal everywhere
unramified abelian extension has Galois group canonically isomorphic
to $\operatorname{Cl}(K)$ by the Artin map.  This isomorphism respects
automorphisms of $K$.  See J.~S.~Milne, \emph{Class Field Theory},
version 4.03, V, Theorem 3.6, Corollary 3.7 and Example 3.9,
pp.~158--159, \url{https://www.jmilne.org/math/CourseNotes/CFT.pdf}.
The equivariance also follows directly by conjugating Frobenius at
unramified primes.

\item[OB-KU: Kummer theory and the unit theorem.]
\label{ob-ref-kummer}
If $K$ contains $\mu_\ell$, then
$H^1(G_K,\mu_\ell)=K^\times/K^{\times\ell}$, with the usual
valuation criterion away from $\ell$.  The group of $S$-units has
rank $r_1+r_2-1+|S|$ and finite cyclic torsion.  See Milne,
\emph{Class Field Theory}, VII, Appendix A, Theorem A.3 and
Proposition A.5, pp.~226--227, and \emph{Algebraic Number Theory},
version 3.08, Theorem 5.11, p.~90,
\url{https://www.jmilne.org/math/CourseNotes/ANT.pdf}.
The Selmer cardinality bound used below is derived from these facts.

\item[OB-MD: ideal representatives and abelian discriminants.]
\label{ob-ref-minkowski}
Every ideal class of a number field of degree $r$ and discriminant
$\Delta$ has an integral representative of norm at most
$(4/\pi)^{r_2}r!r^{-r}\sqrt{|\Delta|}$ (Milne,
\emph{Algebraic Number Theory}, Theorem 4.3, p.~70).
For an abelian field, $|\Delta|$ is the product of the conductors of
its characters (R.~Sharifi, \emph{Iwasawa Theory}, equation (4.2.1),
\url{https://math.ucla.edu/~sharifi/notes/iwasawa-ch04.html}).

\item[OB-BC: Bernoulli denominators.]
\label{ob-ref-bernoulli}
For positive even $j$, von Staudt--Clausen implies that the
denominator of $B_j$ divides $\prod_{p-1\mid j}p$; also
$B_1=-1/2$ and odd $B_j$ vanish for $j>1$.
See T.~Ciurca, \emph{Arithmetic of cyclotomic fields},
Theorem 6.1.5, p.~111,
\url{https://sites.pitt.edu/~caw203/pdfs/Ciurca--Arithmetic_of_cyclotomic_fields--UROP2018.pdf}.

\item[OB-LF: finite extensions of local fields.]
\label{ob-ref-local}
The integers of a finite extension of $\mathbb Q_\ell$ form a finite
free $\mathbb Z_\ell$-module and a complete DVR, with the unique
extended absolute value.  See Milne, \emph{Algebraic Number Theory},
Theorem 7.38, Corollaries 7.39 and 7.42, and Remark 7.43,
pp.~123--124.  The length and norm formulas needed here are proved
explicitly below.
\end{description}

The main project-specific interface is the theorem supplied by OS9.  Let
$\rho:(\mathbb Z/N)^\times\twoheadrightarrow G=\operatorname{Gal}(E/\mathbb Q)$
be the actual cyclotomic restriction map.  For a finite DVR
$O\subset\overline{\mathbb Q}_\ell$, finite over $\mathbb Z_\ell$,
and a character $\psi:G\to O^\times$, put
$\chi=\psi\circ\rho$, extended by zero on nonunits.  For $k\ge2$ put
\[
 P_k=\prod_{p\le k+1}p,\qquad d(N,k)=NP_k,\qquad
 \alpha_{N,k}(\chi)=d(N,k)B_{k,\chi^*},\qquad
 B(N,k)=16N^{k+1}k!,                         \tag{OB.3}
\]
where $\chi^*$ is the primitive inducing character.  These are the
source's \texttt{den}, \texttt{alN}, \texttt{Bprim}, and \texttt{bnd}.
Let $\mathcal S_{O}(E,e)$ be the absolute-Galois version of (OB.1)
with coefficients $(\operatorname{Frac}O/O)(-e)$, and let
$\mathcal S_O(E,e)^\psi$ be its $\psi$-eigenspace for the actual
$G$-action.  This is the source's \texttt{shaEig} in
\texttt{psiLen}; the comparison is proved in OB3.

The hypothesis \textbf{PB} is the literal \texttt{PsiBound}:
for every such original $E,\Omega,N,\rho$, every odd prime
$\ell\nmid n$, every such $O$, every $e\ge1$, every $\psi$ with
$\chi(-1)=(-1)^{e+1}$, and every $a\in O$ whose image is
$\alpha_{N,e+1}(\chi)$,
\[
       \operatorname{length}_O\mathcal S_O(E,e)^\psi
                         \le v_O(a).                \tag{OB.4}
\]
Here $v_O$ is the DVR valuation with value $1$ on a uniformizer and
$v_O(0)=\infty$.  The interface is not replaced by an estimate over
a different field or an unspecified isomorphic coefficient module.

\paragraph{OB1: cyclotomic actions, character count, and the exceptional set.}
\label{claim:OB1}
The given embedding into a cyclotomic field produces a surjective
$\rho$ with
\[
 G\simeq(\mathbb Z/N)^\times/\ker\rho,
 \qquad |G|=[E:\mathbb Q]=n,
 \qquad \rho(-1)=c\ne1.                         \tag{OB.5}
\]
Its lift to a multiple of $N$ is obtained by reduction of units.
For any prescribed parity, the number of complex characters of
$(\mathbb Z/N)^\times$ trivial on $\ker\rho$ is at most $n$.
If $E$ satisfies \texttt{IsCompField}$(M,s,q,d;E)$, with
$o=M.\mathrm{deg}$ and the stated \texttt{PrimePair} hypotheses,
then every odd prime divisor of $n$ divides $o$:
\[
                 \ell\mid n,\quad\ell\ne2
                       \quad\Longrightarrow\quad\ell\mid o.
 \tag{OB.6}
\]

\begin{proof}
The cyclotomic Galois group is $(\mathbb Z/N)^\times$.
All of its subgroups are normal, so the image of the given embedding
of $E$ is Galois over $\mathbb Q$, and restriction is surjective.
Choose an embedding of this cyclotomic overfield into the already
given $\Omega$ extending the original embedding of $E$.  An
automorphism of $\Omega$ acting as the $b$th power on $\mu_N$
therefore restricts to this very $\rho(b)$ on $E$.  This is
\texttt{IsCycAct}.  Restriction at a larger modulus commutes with
reduction of units.  Complex conjugation sends every root of unity
to its inverse, giving (OB.5); its restriction is nontrivial because
$E$ is CM.  Distinct characters of a finite abelian group are linearly
independent as functions, by the usual orthogonality sum, so any
subset of the character group has size at most $|G|$.

For (OB.6), work in the actual cyclotomic overfield supplied by
\texttt{IsCompField}.  Its fixing condition says that an exponent
$b$ fixes $E$ precisely when it is $\pm1$ modulo $2^{s+2}$,
belongs to the kernel specifying the abelian base $M$, and has
Jacobi symbol $1$ at $q$ and $d$.  Every cyclotomic automorphism
raised to $R=2^{s+1}o$ meets all these conditions: Euler's theorem
handles $2^{s+2}$, the finite base quotient has order $o$, and the
even exponent makes both Jacobi symbols $1$.  Thus the exponent of
$\operatorname{Gal}(E/\mathbb Q)$ divides $R$.  Cauchy's theorem
gives an element of order $\ell$ if $\ell\mid n$, so
$\ell\mid2^{s+1}o$; an odd $\ell$ must divide $o$.
This argument uses the original fixing subgroup and degree.
\end{proof}

\paragraph{OB2: the integral Bernoulli product and its prime valuation budget.}
\label{claim:OB2}
For $k\ge2$, fix a set $P$ of subgroups of $(\mathbb Z/N)^\times$
and let $\mathcal C$ consist of all characters with $\ker\chi\in P$
and $\chi(-1)=(-1)^k$.  There is a nonzero integer $J_{N,k,\mathcal C}$
such that, in every characteristic-zero splitting field,
\[
 \prod_{\chi\in\mathcal C}\alpha_{N,k}(\chi)
     =J_{N,k,\mathcal C},\qquad
 \log|J_{N,k,\mathcal C}|\le|\mathcal C|\log B(N,k).
 \tag{OB.7}
\]
Put $w_\ell=v_\ell(J_{N,k,\mathcal C})$.  For any finite set
$T_\ell$ of the corresponding $\ell$-adic characters and any finite
set $S$ of primes,
\[
 \sum_{\chi\in T_\ell}-\log|\alpha_{N,k}(\chi)|_\ell
       \le w_\ell\log\ell,\qquad
 \sum_{\ell\in S}w_\ell\log\ell
       \le |\mathcal C|\log B(N,k).                 \tag{OB.8}
\]

\begin{proof}
For a character $\eta$ of modulus $f$, the generalized Bernoulli
number is
\[
 B_{k,\eta}=f^{k-1}\sum_{a=1}^{f}\eta(a)B_k(a/f),\qquad
 B_k(X)=\sum_{j=0}^k\binom kj B_jX^{k-j}.          \tag{OB.9}
\]
OB-BC shows that $P_kB_j$ is integral for $0\le j\le k$.
Multiplication of (OB.9) by $fP_k$ consequently yields an integer
linear combination of the values of $\eta$.  If $f\mid N$, the
same is true after multiplication by $NP_k$.  The values of the
primitive inducer occur among those of the original character:
reduction $(\mathbb Z/N)^\times\to(\mathbb Z/f)^\times$ is
surjective, as follows by choosing lifts avoiding the additional
prime divisors of $N$.  Therefore $\alpha_{N,k}(\chi)$ belongs
to the ring generated by the values of $\chi$, and is an algebraic
integer.  In particular it lies in the actual $O$ whenever those
values lie in $O$.

Here are elementary bounds with the source's constant.  Integration
by parts in the Fourier coefficients of $B_k$, using
$B_k'=kB_{k-1}$ and $B_k(1)=B_k(0)$ for $k\ge2$, gives
\[
 B_k(x)=-\frac{k!}{(2\pi i)^k}
             \sum_{r\ne0}\frac{e^{2\pi irx}}{r^k}.
\]
For $k\ge2$ this converges absolutely and
$\zeta(k)\le1+\int_1^\infty u^{-k}du\le2$.
Thus $|B_{k,\eta}|\le4f^kk!/(2\pi)^k$.
The elementary inequality $\prod_{p\le h}p\le4^h$ follows by
induction: primes in $(m,2m]$ divide $\binom{2m}{m}$; primes in
$(m+1,2m+1]$ divide $\binom{2m+1}{m}\le4^m$, since the two
central binomial coefficients are equal.  Applying the induction
hypothesis to $m$ or $m+1$ proves the claim.  Since $4<2\pi$,
\[
 |\alpha_{N,k}(\chi)|
 \le N4^{k+1}\frac{4N^kk!}{(2\pi)^k}
 \le16N^{k+1}k!=B(N,k).                          \tag{OB.10}
\]

Nonvanishing here requires no result at $s=1$.  For primitive
$\eta$ of parity $(-1)^k$, insert the Fourier series into (OB.9).
The finite Gauss-sum identity gives
\[
 B_{k,\eta}
   =-\frac{2k!f^{k-1}\tau(\eta)}{(2\pi i)^k}
                              L(k,\overline\eta).
\]
The Gauss sum is nonzero: character orthogonality gives
$|\tau(\eta)|^2=f$; for the trivial conductor-one character its
value is $1$.  The absolutely convergent Euler product at $k>1$
is nonzero.  One can verify the Gauss identity for all integer
frequencies by substitution when the frequency is coprime to $f$;
otherwise primitivity supplies a unit that fixes the additive
character and has nontrivial multiplicative value, forcing its
sum to vanish.  Parseval on the finite cyclic group then gives
the stated squared norm.  Thus every factor in (OB.7) is nonzero.

Place the character values in one finite cyclotomic field.
Its Galois automorphisms raise these values to powers coprime to
their order, preserving their kernels and the parity.  They
therefore permute the factors of (OB.7).  The product is rational
and integral, hence an integer; it is nonzero by the preceding
paragraph.  The construction uses rational polynomials in the
roots of unity, so the same integer is obtained in every splitting
field.  Taking the product of (OB.10) gives its logarithmic bound.

In an $\ell$-adic splitting field each factor is integral, so its
absolute value is at most $1$.  The negative logarithms are
nonnegative, and their sum over the full set is
$-\log|J|_\ell=v_\ell(J)\log\ell$.  A partial sum is no larger.
Unique factorization of the nonzero integer $J$ gives
$\sum_{\ell\in S}v_\ell(J)\log\ell\le\log|J|$.
These prove (OB.8), including finiteness of the prime budget.
\end{proof}

\paragraph{OB3: the tame-prime length bound on the original local data.}
\label{claim:OB3}
Assume PB.  For every odd prime $\ell\nmid n$ and every $e\ge1$,
$m_\ell(E,e;D)$ is finite.  With
\[
 \mathcal C=\{\chi:\ker\rho\subseteq\ker\chi,
                          \ \chi(-1)=(-1)^{e+1}\},
\]
one has, for every finite set $S$ of odd primes not dividing $n$,
\[
 \sum_{\ell\in S}\log\ell\,m_\ell(E,e;D_\ell)
                \le n\log B(N,e+1).                 \tag{OB.11}
\]
The local-data system $D_\ell$ is arbitrary, and the same $E$ and
$\Omega$ occur in all terms.

\begin{proof}
First identify (OB.1) with the strict Selmer subgroup in absolute
Galois cohomology: its classes vanish on decomposition groups above
$\ell$ and on all inertia groups.  Away from $\ell$, inertia acts
trivially on $A_\ell(e)$; vanishing of its cohomology class is
therefore literal vanishing of a representative cocycle there.
Such a cocycle vanishes on the closed normal subgroup generated
by these inertia groups, so factors through $G_{E,S}$.
Conversely an inflated cocycle is unramified there.  The conditions
at $\ell$ are exactly the local coboundary conditions in (OB.1).
Changing the chosen place above a prime conjugates its decomposition
group and preserves this condition.  This proves the comparison
for the actual $D$, without choosing another host field.

Choose a finite extension of $\mathbb Q_\ell$ inside
$\overline{\mathbb Q}_\ell$ containing the $|G|$th roots of unity,
and let $O$ be its integer ring.  Write
$h=v_O(\ell)=\operatorname{length}_O(O/\ell O)$; this is a positive
finite integer.  Since $O$ is finite free over $\mathbb Z_\ell$,
tensoring continuous cochains with $O$ is the same as taking a
finite direct sum of their coordinate cochains.  It commutes with
cocycles, coboundaries and all local restriction kernels, including
the full family of inertia conditions.  Also
$O\otimes_{\mathbb Z_\ell}(\mathbb Q_\ell/\mathbb Z_\ell)
=\operatorname{Frac}(O)/O$.  Hence
\[
 O\otimes\mathcal S_\ell(E,e;D)=\mathcal S_O(E,e),
 \quad
 \operatorname{length}_O(O\otimes X)
             =h\operatorname{length}_{\mathbb Z_\ell}X.
 \tag{OB.12}
\]
For finite length the second formula follows by tensoring a
composition series and using the length of $O/\ell O$.
For arbitrary torsion modules it follows by taking the supremum
over finitely generated submodules.  All cohomology classes here
are $\ell$-power torsion, since a continuous cocycle from a compact
group into the discrete torsion module has finite image.

Conjugation and coefficient action commute with the comparison in
(OB.12).  Inner automorphisms act trivially on cohomology, as the
cocycle identity expresses their difference from the identity by
a coboundary.  Thus the action factors through the original
$G=\operatorname{Gal}(E/\mathbb Q)$, with exactly the sign in (OB.2).
Since $\ell\nmid |G|$, the elements
\[
 e_\psi=|G|^{-1}\sum_{g\in G}\psi(g)^{-1}[g]\in O[G]
\]
are mutually orthogonal projectors summing to $1$.
This follows by summing finite characters on $G$; the character
set has $|G|$ elements, by decomposing $G$ into cyclic factors.
The $(-1)^{e+1}$ eigenspace of $c$ is therefore the direct sum of
$\mathcal S_O(E,e)^\psi$ for
$\psi(c)=(-1)^{e+1}$.  These are exactly the \texttt{psiLen}
objects in PB, obtained from \texttt{shaEig} by factoring out the
inner action, not different cohomology groups selected by dimension.

OB2 supplies $a_\psi\in O$ representing
$\alpha_{N,e+1}(\psi\circ\rho)$ and proves $a_\psi\ne0$.
PB and (OB.12) give
\[
 h\,m_\ell(E,e;D)
  =\sum_{\psi(c)=(-1)^{e+1}}
                  \operatorname{length}_O\mathcal S_O(E,e)^\psi
  \le\sum_{\psi(c)=(-1)^{e+1}}v_O(a_\psi)<\infty.
 \tag{OB.13}
\]
Writing $\ell=u\varpi^h$ and $a=v\varpi^{v_O(a)}$ shows
\[
            v_O(a)\log\ell=-h\log|a|_\ell.
\]
Indeed units have absolute value $1$, because both the unit and
its inverse are integral, and $|\ell|_\ell=\ell^{-1}$.
Multiplying (OB.13) by $\log\ell$ and canceling the positive $h$
gives a sum of the logarithms in (OB.8).  Surjectivity of $\rho$
makes $\psi\mapsto\psi\circ\rho$ injective, so no character is
counted twice.  Finally OB2 and $|\mathcal C|\le n$ from OB1
prove (OB.11).  Its integer product depends on $E,N,e$, not on
the varying prime or the choices of $D_\ell$ and $O$.
\end{proof}

\paragraph{OB4: a uniform bound on the number of cohomology generators.}
\label{claim:OB4}
For any number field $E\hookrightarrow\mathbb Q(\zeta_N)$,
any prime $\ell$, any integer $e$, and any original $D$, the
$\ell$-torsion of $H^1(G_{E,S},A_\ell(e))$ has finite length $g$ and
\[
                    g\le8\ell n\log(\ell N).
 \tag{OB.14}
\]

\begin{proof}
Set $K=E(\mu_\ell)$ inside the given $\Omega$.  Then
$[K:E]\mid\ell-1$, $K$ contains $\mu_\ell$, and the given
cyclotomic embedding extends to
$K\hookrightarrow\mathbb Q(\zeta_{\ell N})$.
Put $r=[K:\mathbb Q]\le\ell n$.
An $\ell$-torsion class has a representative cocycle taking values
in $A_\ell(e)[\ell]$: if $\ell f=\delta b$, choose $b'$ with
$\ell b'=b$ and replace $f$ by $f-\delta b'$.
The action on these values becomes trivial over $K$, so restriction
is an additive character $G_K\to\mathbb F_\ell$, unramified
outside the primes over $\ell$.

Choosing one such representative for each class gives an injective
map into this set of characters.  To see injectivity, if two
restrictions agree, the difference cocycle vanishes on $G_K$.
It descends to $\operatorname{Gal}(K/E)$ with values in the
$G_K$-invariants.  Its order is prime to $\ell$; averaging a
cocycle over this finite group writes it as a coboundary, since
division by the group order is available on every
$\mathbb Z_\ell$-module.  The original two classes are equal.
This argument asserts a set injection, and does not incorrectly
assume that the chosen representatives define a canonical linear map.

By Kummer theory, the target injects into
\[
 \operatorname{Sel}_S(K,\ell)
 =\{a\in K^\times:v_{\mathfrak p}(a)\in\ell\mathbb Z
                  \ (\mathfrak p\notin S)\}/K^{\times\ell}.
\]
The valuation condition follows by taking the valuation of
$\beta^\ell=a$ in an unramified local extension away from $\ell$.
Factoring the principal ideal outside $S$ gives the exact sequence
\[
 0\longrightarrow O_{K,S}^\times/(O_{K,S}^\times)^\ell
 \longrightarrow\operatorname{Sel}_S(K,\ell)
 \longrightarrow\operatorname{Cl}_S(K)[\ell]
 \longrightarrow0.
 \tag{OB.15}
\]
The last map sends $a$ to the class of the ideal whose $\ell$th
power is $(a)$ away from $S$.  Its kernel and surjectivity are
immediate from this construction.  There are at most $r$ primes
over $\ell$.  OB-KU therefore bounds the first term in (OB.15)
by $\ell^{2r}$, and the second class-group term by $h_K$.
Consequently the original $\ell$-torsion group has at most
$\ell^{2r}h_K$ elements.  Being killed by $\ell$, it is an
$\mathbb F_\ell$-vector space of dimension and length $g$, so
\[
                    g\log\ell\le2r\log\ell+\log h_K.
 \tag{OB.16}
\]

We record the class-number bound with its constants.  If
$K\hookrightarrow\mathbb Q(\zeta_L)$ and $L\ge2$, OB-MD gives
$|\Delta_K|\le L^r$ and $r\le\varphi(L)\le L$.
Minkowski's constant is at most $1$: for $r\ge2$, use
$r!\le((r+1)/2)^r$ and $(r+1)/(2r)\le3/4$; the degree-one
case is immediate.  Each class has a representative of norm at
most $X=\sqrt{|\Delta_K|}$.
The number of ideals of norm at most $X$ is at most
$X(1+\log X)^{r-1}$: at every rational prime the ideal-counting
Euler factor is coefficientwise at most $(1-T)^{-r}$, and the
sum of the $r$-fold divisor function up to $X$ is bounded by
$X(1+\log X)^{r-1}$ by induction using the harmonic sum.
It follows that
\[
 \log h_K\le\tfrac12\log|\Delta_K|
       +(r-1)\log(1+\tfrac12\log|\Delta_K|)
                       \le3r\log L.                 \tag{OB.17}
\]
For the last inequality use $\log|\Delta_K|\le r\log L$ and
$1+\tfrac12r\log L\le L^2$, which follows from
$r\le L$ and $\log L\le L-1$.
Apply (OB.17) with $L=\ell N$.  Then (OB.16) is at most
$5r\log(\ell N)$.  Since $\log\ell\ge\log2>5/8$,
$g\le8r\log(\ell N)\le8\ell n\log(\ell N)$, as claimed.
\end{proof}

\paragraph{OB5: an actual annihilator of the negatively twisted minus group.}
\label{claim:OB5}
For an odd prime $\ell$ and $e\ge1$, put $R=N\ell$,
$a=2R+1$, $k=e+1$, and let
$\pi:(\mathbb Z/R)^\times\twoheadrightarrow G$
be the lift of the original restriction map.  In $\mathbb Q[G]$
define
\[
 \begin{split}
 E_{L,k}(b)&=\frac{L^{k-1}}{k}B_k(b/L),\\
 y&=(1-a^k[\pi(a)])
                  \left(-\sum_{b\in(\mathbb Z/R)^\times}
                                  E_{R,k}(b)[\pi(b)]\right).
 \end{split}                                                   \tag{OB.18}
\]
Then $y$ has coefficients in $\mathbb Z_\ell$, and its resulting
element $z\in\mathbb Z_\ell[G]$ satisfies
\[
                 z\,\varphi(R)x=0
                 \qquad(x\in\mathcal S_\ell(E,e;D)).
 \tag{OB.19}
\]
The source writes this element as
\texttt{yEl (e+1) a (piQ rho)\^{}(-1)}: the inverse in that
argument compensates for the inverse used in \texttt{sumEl},
leaving exactly the group elements $[\pi(b)]$ in (OB.18).

\begin{proof}
For $L=R\ell^m$, let $b$ denote a unit representative in $[0,L)$,
let $b'$ represent $a^{-1}b$ modulo $L$, and set
\[
 j_L(b)=\frac{ab'-b}{L}\in\mathbb Z,\qquad
 E_1(L,a,b)=\frac{a-1}{2}-j_L(b)\in\mathbb Z,
 \qquad
 z_m=-\sum_{b\in(\mathbb Z/L)^\times}
                  b^e E_1(L,a,b)[\pi(b)].             \tag{OB.20}
\]
Here $a$ is odd and coprime to every $L$.
We first prove $z_m\to y$ coefficientwise $\ell$-adically.
The first two terms of the Bernoulli polynomial give
\[
 E_{L,k}(b)=\frac{b^k}{kL}-\frac{b^{k-1}}2
                      +\frac{L}{kP_k}u_b,
                  \qquad u_b\in\mathbb Z.
\]
Indeed all remaining terms have a factor $L$ after their
denominators are multiplied by $kP_k$.
Since $ab'=b+Lj_L(b)$, the binomial theorem, first modulo $L^2$
for the $k$th powers and then modulo $L$ for the $(k-1)$st powers,
shows
\[
 E_{L,k}(b)-a^kE_{L,k}(b')
      =b^eE_1(L,a,b)+\frac{L}{2kP_k}v_b,
                         \qquad v_b\in\mathbb Z.
 \tag{OB.21}
\]
Thus the error has arbitrarily large $\ell$-adic valuation as
$m\to\infty$, uniformly over its finitely many coefficients.
The distribution identity
$\sum_{j=0}^{h-1}B_k((x+j)/h)=h^{1-k}B_k(x)$
follows at once by summing the defining generating function.
Since $\ell\mid R$, every lift of a unit modulo $R$ to $R\ell^m$
is a unit.  This identity proves that the higher-level element in
(OB.18), projected to $G$, is the same $y$.  Reindexing multiplication
by $a$ gives the coefficient difference in (OB.21).  Summing
(OB.21) therefore proves that $y-z_m$ belongs to
$\ell^r\mathbb Z_\ell[G]$ once
$m\ge r+v_\ell(2kP_k)$.  In particular $y$ is integral at $\ell$;
its $\ell$-adic realization is $z$.

We now prove the cohomological annihilation, including the twist.
Pull a cocycle $f$ representing $x\in\mathcal S_\ell(E,e;D)$
back to $G_E$.  Its image is finite, hence killed by some
$\ell^u$.  At each of the finitely many primes of $E$ over
$\ell$, its restriction is $\delta b_v$ for an element $b_v$
of $A_\ell(e)$.  Increase $u$ so that all $b_v$ are killed by
$\ell^u$.  For $m$ large enough, cyclotomic values on
$G_{K_m}$, where $K_m=\mathbb Q(\zeta_{R\ell^m})\subset\Omega$,
are $1$ modulo $\ell^u$.  Consequently
$f|_{G_{K_m}}$ is an additive character and vanishes on the
decomposition groups over $\ell$.  It is unramified at every
other finite prime because the original cocycle factors through
$G_{E,S}$.  These assertions hold at every conjugate place:
the cocycle conjugation identity adds only a term
$(1-\chi_{\rm cyc}(g)^{-e})f(h)$, which is zero on this level.
Since $K_m$ is totally complex, OB-CF makes this restriction an
ordinary class-group character.

We use the finite Jacobi annihilator proved in OS1, keeping its
group-ring orientation.  Increase $m$ so that $a<L$.  For a unit
$t$ modulo $L$, represented in $(0,L)$, its coefficient there is
\[
 \theta_{L,a}(t)
 =\left\lfloor\frac{at}{L}\right\rfloor
       -\left\lfloor\frac{a(L-t)}{L}\right\rfloor
 =2\left\lfloor\frac{at}{L}\right\rfloor-(a-1).
\]
The last equality holds because $at/L$ is not an integer.  Put
$b\equiv at\pmod L$, so $b'=t$ and
$j_L(b)=\lfloor at/L\rfloor$.  Reindexing gives the exact identity
\[
 2\sum_bE_1(L,a,b)\sigma_b^{-1}
       =-\sigma_a^{-1}
               \sum_t\theta_{L,a}(t)\sigma_t^{-1}.
\]
OS1 constructs a principal ideal for the right-hand sum applied to
every ideal, so that sum annihilates $\operatorname{Cl}(K_m)$.
The values of the class-group character $f|_{G_{K_m}}$ have
odd $\ell$-power order; multiplication by $2$ is invertible on
them.  Artin equivariance therefore gives, for $v\in G_{K_m}$
and lifts $\sigma_b$ of the cyclotomic automorphisms,
\[
              \sum_b E_1(L,a,b)f(\sigma_b^{-1}v\sigma_b)=0.
 \tag{OB.22}
\]
For comparison only, the source's classical Stickelberger route is
the integral identity
$\sum_bj_L(b)\sigma_b^{-1}=-\sigma_a^{-1}(\sigma_a-a)\theta_L$,
with $\theta_L=L^{-1}\sum_b b\sigma_b^{-1}$.
Its classical annihilation theorem is Sharifi,
\emph{Iwasawa Theory}, Definition 4.5.1, Lemma 4.5.2 and
Theorem 4.5.6.  It is not a retained input: the OS1 argument just
given proves exactly the odd-primary consequence needed here.

Let $U=G_{E(\mu_{\ell^\infty})}$ and let
$v\in U\cap G_{\mathbb Q(\mu_R)}$.  Such $v$ belongs to every
$G_{K_m}$.  On restriction to $U$, the action on a cohomology class is
\[
 \operatorname{res}(\pi(b)x)(v)
   =\chi_{\rm cyc}(\sigma_b)^{-e}
                       f(\sigma_b^{-1}v\sigma_b).
 \tag{OB.23}
\]
The lift satisfies $\chi_{\rm cyc}(\sigma_b)\equiv b\pmod{\ell^m}$.
Because its values are units and the displayed cocycle values are
killed by $\ell^u$, multiplication by
$b^e\chi_{\rm cyc}(\sigma_b)^{-e}$ acts as $1$ for $m\ge u$.
Equations (OB.20), (OB.22), and (OB.23) now give
$\operatorname{res}(z_mx)(v)=0$.  The coefficient congruence
$z-z_m\in\ell^u\mathbb Z_\ell[G]$ allows $z_m$ to be replaced
by $z$.  This explicitly accounts for the negative twist.

Finally restriction to $U$ is injective on the original $H^1$.
Indeed its kernel is represented by cocycles descending to
$\operatorname{Gal}(E(\mu_{\ell^\infty})/E)$, an infinite open
subgroup of $\mathbb Z_\ell^\times$.  For odd $\ell$ it is
procyclic: its torsion part has order prime to $\ell$ and its
principal-unit part is a copy of $\mathbb Z_\ell$.
If $\gamma$ is a topological generator, then
$\chi_{\rm cyc}(\gamma)^{-e}-1\ne0$, because $e\ge1$ and
$\gamma$ has infinite order.  Multiplication by this nonzero
$\ell$-adic integer is surjective on
$\mathbb Q_\ell/\mathbb Z_\ell$.  Subtract a coboundary agreeing
with the cocycle at $\gamma$; the remaining cocycle is zero on
the dense cyclic subgroup and hence everywhere by continuity.
This proves injectivity.  The homomorphism
$\operatorname{res}(zx):U\to A_\ell(e)$ vanishes on
$U\cap G_{\mathbb Q(\mu_R)}$.  Its quotient is a subgroup of
$(\mathbb Z/R)^\times$, so its exponent divides $\varphi(R)$.
Thus $\varphi(R)\operatorname{res}(zx)=0$, and injectivity gives
(OB.19).  No infinite-level finiteness assertion or implicit change
of the local-data system has been used.
\end{proof}

\paragraph{OB6: a determinant gives an explicit annihilating power.}
\label{claim:OB6}
For every odd $\ell$, every $e\ge1$, and every original $D$, there
is $r\ge0$ such that
\[
 \ell^r\mathcal S_\ell^-(E,e;D)=0,\qquad
      r\log\ell\le30n(e+1)\log(e\ell N).
 \tag{OB.24}
\]

\begin{proof}
Keep $R,a,k,G,z,y$ from OB5 and put $\varepsilon=(-1)^k$.
Because $c\in G$ is a nonidentity involution,
\[
       Q=\mathbb Z_\ell[G]/([c]-\varepsilon)
\]
is free of rank $n/2$: choose one representative from each pair
$\{g,cg\}$, replacing the other by $\varepsilon g$.
Let $\delta$ be the determinant of multiplication by $z$ on $Q$.
The adjugate identity, applied to $1\in Q$, gives
$\delta\in(z,[c]-\varepsilon)$ in $\mathbb Z_\ell[G]$.
Together with (OB.19) and $cx=\varepsilon x$, it follows that
$\delta\varphi(R)x=0$ for every $x$ in the minus group.

Put
\[
 D_{R,k}=Rk(k+1)!,\quad
 S_{R,a,k}=2a^kR^kk!,\quad
 X_{R,a,k}=D_{R,k}S_{R,a,k}.                    \tag{OB.25}
\]
Multiplying $y$ by $D_{R,k}$ gives an integral group-ring element:
the Bernoulli polynomial expansion and OB-BC clear every coefficient.
The determinant on the integral signed quotient is therefore a
rational integer $J$, and
\[
      J=D_{R,k}^{\,n/2}\delta,\qquad
      0<|J|\le X_{R,a,k}^{\,n/2}.                  \tag{OB.26}
\]
Here is a proof of nonvanishing and the bound.  Over $\mathbb C$ the
signed quotient diagonalizes into the $n/2$ characters $\eta$ of
$G$ with $\eta(c)=\varepsilon$.  The eigenvalue of $y$ is
\[
       (1-a^k\chi(a))\frac{-B_{k,\chi}}{k},
               \qquad \chi=\eta\circ\pi.
 \tag{OB.27}
\]
The smoothing factor is nonzero, since $a>1$ and $|\chi(a)|=1$.
For an imprimitive character of level $R$, the distribution formula
in (OB.9), excluding the extra nonunits prime by prime, gives
\[
 B_{k,\chi}=B_{k,\chi^*}
       \prod_{p\mid R,\ p\nmid\operatorname{cond}\chi}
                         (1-\chi^*(p)p^{k-1}).
\]
Each extra factor is nonzero for $k\ge2$, and the primitive factor
is nonzero by OB2.  Thus (OB.27) is nonzero.  The same Fourier bound
as in OB2 gives $|B_{k,\chi}|\le R^kk!$, because
$4/(2\pi)^k\le1$.  Also
$|1-a^k\chi(a)|\le2a^k$ and $k\ge1$.
This proves the bound $S_{R,a,k}$ for each eigenvalue, hence (OB.26).

Set $r_0=v_\ell(J)$ and $r_1=v_\ell(\varphi(R))$.
Since $\delta\mid J$ in $\mathbb Z_\ell$, it divides
$\ell^{r_0}$ there; similarly $\varphi(R)\mid\ell^{r_1}$.
Thus $r=r_0+r_1$ kills the minus group.  Moreover
\[
 r\log\ell\le n\log X_{R,a,k}+\log R
                    \le n\log(X_{R,a,k}R).          \tag{OB.28}
\]
The first inequality even permits $n/2$ in place of $n$ in its
first term; the weaker form is the one used by the source.

For completeness, its numerical constant is explicit.  Put
$U=e\ell N\ge2$ and $T=3U$.  Then
$R\le T$, $a=2R+1\le T^3$, $k=e+1\le T$,
$k+1\le T^2$, and $k\ge2$.  Hence
\[
 \begin{split}
 X_{R,a,k}R
 &=Rk(k+1)!\,2a^kR^kk!\,R\\
 &\le T^{7k+6}\le T^{10k}.
 \end{split}
\]
We used $(k+1)!\le(T^2)^{k+1}$ and $k!\le T^k$.
Since $3\le U^2$, $\log T\le3\log U$.
Substituting in (OB.28) proves (OB.24) with $30$.
The construction fixes $r$ before $x$ varies; it is not a
pointwise collection of unrelated annihilating exponents.
\end{proof}

\paragraph{OB7: exponent times generators gives the exceptional-prime bound.}
\label{claim:OB7}
For every odd prime $\ell$, every $e\ge1$, and every $D$,
\[
 m_\ell(E,e;D)<\infty,\qquad
 m_\ell(E,e;D)\log\ell
       \le240\ell n^2(e+1)\log(e\ell N)^2.
 \tag{OB.29}
\]

\begin{proof}
For an endomorphism $f$ of any module, the map induced by $f$ from
$\ker f^{j+1}$ to $\ker f^j$ has kernel $\ker f$.
The length inequality for this exact sequence proves inductively
that $\operatorname{length}\ker f^r\le
r\operatorname{length}\ker f$.
Apply this to multiplication by $\ell$ on the ambient $H^1$.
OB6 puts the minus group inside $\ker\ell^r$, while OB4 bounds
the length $g$ of $\ker\ell$.  Thus its length is finite and at
most $rg$.  Multiplication by $\log\ell$ and (OB.14), (OB.24)
give
\[
 m_\ell\log\ell
   \le240\ell n^2(e+1)\log(e\ell N)\log(\ell N)
   \le240\ell n^2(e+1)\log(e\ell N)^2.
\]
Both logarithms are nonnegative and $e\ge1$.  This is precisely
the odd-prime portion of \texttt{StickBound 240} needed below.
No claim about its unused $\ell=2$ extension is required.
\end{proof}

\paragraph{OB8: the full uniform \texttt{ThmBcFree} supplier.}
\label{claim:OB8}
Use PB from OS9, the Jacobi annihilator from OS1, and the explicitly
listed classical theorems.  Let
$M=(m,H)$ be any \texttt{AbF}, put $o=[(\mathbb Z/m)^\times:H]$,
let $s,t\ge2$, and let $q,d$ satisfy \texttt{PrimePair}$(m,q,d)$.
For every original CM number field $E$ satisfying
\texttt{IsCompField}$(M,s,q,d;E)$, put
$n=[E:\mathbb Q]$ and $F=\operatorname{condF}(E)$.
Then for every given algebraic closure $\Omega$, every family of
local data $(D_\ell)$ on that same $E,\Omega$, and every finite
set $S$ of odd primes,
\[
 \begin{split}
 &m_\ell(E,e;D_\ell)<\infty
                 \qquad(\ell\text{ odd},\ 1\le e<t),\\
 &\sum_{\ell\in S}\log\ell\sum_{e=1}^{t-1}
                         m_\ell(E,e;D_\ell)\\
 &\qquad\le3nt^2\log(tF)
       +240\sum_{\substack{\ell\mid o\\\ell\text{ odd prime}}}
                       \ell n^2t^2\log(t\ell F)^2.
 \end{split}                                                   \tag{OB.30}
\]
Thus the literal \texttt{ThmBcFree} holds with $C=240$, and it
supplies VW.13--VW.17 with that same constant.

\begin{proof}
The set defining $F$ is nonempty by \texttt{IsCompField}.  Its
least element is positive and itself supplies an embedding of the
original $E$ into $\mathbb Q(\zeta_F)$.  Apply OB3 and OB7 with
this $N=F$.  They prove all asserted individual finiteness
statements before any real-valued lengths are summed.

Split $S$ into $S_0=\{\ell\in S:\ell\nmid n\}$ and its complement.
For $1\le e<t$, (OB.11) bounds the sum over $S_0$ by
$n\log B(F,e+1)$.  The source's elementary numerical bound is
\[
        B(F,e+1)=16F^{e+2}(e+1)!
                       \le(tF)^{2t+5}.              \tag{OB.31}
\]
Indeed $16\le(tF)^4$, $F^{e+2}\le(tF)^{t+1}$, and
$(e+1)!\le t^t\le(tF)^t$.  Interchanging the two finite sums and
using $(t-1)(2t+5)\le3t^2$ for integers $t\ge2$ gives
\[
       \sum_{\ell\in S_0}\log\ell
                    \sum_{e=1}^{t-1}m_\ell
                        \le3nt^2\log(tF).           \tag{OB.32}
\]
For a prime in the complementary set, (OB.29) and
$e+1\le t$, $\log(e\ell F)\le\log(t\ell F)$ give
\[
      \log\ell\sum_{e=1}^{t-1}m_\ell
              \le240\ell n^2t^2\log(t\ell F)^2.
 \tag{OB.33}
\]
All summands are nonnegative.  OB1 puts every prime in this
complementary set among the odd prime divisors of the fixed
base degree $o$, so (OB.33) may be summed over that full set.
Combining it with (OB.32) gives (OB.30).  The same calculation works
for every $S$, with no need to assert beforehand that only finitely
many prime contributions are nonzero.  The resulting right side
is exactly \texttt{minusBound 240 o n F t}, and the universal
quantifiers are exactly those of \texttt{VVminusLe} and
\texttt{ThmBcFree}.
\end{proof}

\paragraph{Dependency consequence.}
OB1--OB8 close the full odd-minus supplier from the literal PB
interface proved in OS9, the Jacobi annihilation proved in OS1,
and the named classical theorems.  Together these replace
\texttt{ThmBcFree} in the structural boundary of VW4; no additional
bespoke odd-minus arithmetic or quantitative interface remains
in this branch.  The resulting uniform
constant is independent of the abelian base, conductor, host primes,
height, twist range, algebraic closure and local-data choices.
This statement leaves the other VW4 suppliers, and the later
lattice-gain/character-balance inputs, as separate obligations.
